% Restored and adapted detailed content from the previous conclusion draft.

% Source: paper/chapters/07-conclusion.tex

This thesis presented the design, implementation, and evaluation of a novel Domain-Specific Language
(domain-specific language) for musical
composition.
By treating musical structures as first-class programming constructs and shifting the complexity of temporal resolution
entirely to a compile-time lowering engine, the project established a new paradigm for offline, deterministic
algorithmic composition.

\section{Summary of Contributions}\label{detail-sec:summary-of-contributions}

The primary contribution of this research is the successful demonstration of a ``zero-runtime'' architecture for a
musical domain-specific language\@.
Unlike live-coding environments that rely on threaded execution and real-time interpreters, this compiler successfully
maps high-level musical intent into a static, consistently quantized Intermediate Representation (IR).

Key achievements include:

\noindent \textbf{Structural Abstraction.} The language provides robust constructs (\texttt{pattern}, \texttt{track},
\texttt{voice}) that allow composers to build complex polyphonic arrangements modularly, without manually
calculating absolute timestamps or delta-times.

\noindent \textbf{Deterministic Lowering.} The implementation of a stateful cursor mutation algorithm successfully
``unrolls'' C-style imperative control flow (loops, conditionals) and recursive pattern calls into a linear sequence
of MIDI events, while consistently preserving temporal gaps and rests.

\noindent \textbf{Type-Safe Semantics.} The custom semantic analyzer rigorously enforces musical logic, ensuring that
operators are applied correctly and that patterns are invoked with valid arguments, thereby preventing invalid data
from reaching the backend.

\noindent \textbf{Robust Tooling.} The development of a full-stack, web-based playground with real-time diagnostics and
Piano Roll visualization validates the compiler's performance and provides an accessible environment for users.

\section{Limitations}\label{detail-sec:limitations}

While the compiler achieves its primary goals, the architectural choices introduce certain limitations:

\noindent \textbf{Lack of Real-Time Interaction.} Because the output is fully resolved at compile-time into a
static file,
the language cannot respond to external input (such as live MIDI controllers or network signals) during playback.

\noindent \textbf{MIDI Expressiveness Constraints.} The current implementation is bound by the constraints of
the MIDI 1.0
standard.
Features like continuous parameter automation (e.g., filter sweeps, pitch bends) or microtonal tunings are not
currently supported by the AST\@.

\noindent \textbf{Static Pattern Signatures.} Pattern parameters currently only accept concrete evaluated types.
The language lacks support for passing un-evaluated functions or lambdas, which limits the potential for
higher-order musical transformations.

\section{Future Work}\label{detail-sec:future-work}

The foundation laid by this compiler opens several promising avenues for future research and development:

\subsection{Standard Library and Advanced Musical
Constructs}\label{detail-subsec:standard-library-and-advanced-musical-constructs}
Future iterations could include a standard library of musical utilities.
Built-in functions for generating scales, arpeggios, and chord inversions would greatly enhance the language's
expressiveness.
Additionally, extending the type system to understand key signatures and functional harmony (e.g., automatically
resolving Roman numeral chord progressions) would elevate Motivo to a higher level of musical abstraction.

\subsection{Continuous Parameter Automation}\label{detail-subsec:continuous-parameter-automation}
To support more expressive performances, the AST and lowering engine could be extended to support continuous data.
Constructs could be added to define ``envelopes'' or ``curves'' over time, which the backend would then serialize into a
dense stream of MIDI Control Change (CC) or Pitch Bend messages.

\subsection{Alternative Backends}\label{detail-subsec:alternative-backends}
The separation of the IR from the backend makes it feasible to support output formats beyond MIDI\@.
A compelling future project would be the development of an Open Sound Control (OSC) backend or a direct audio synthesis
backend (e.g., generating Csound or SuperCollider scripts), allowing the language to directly interface with modern
digital signal processing environments.

\vspace{1cm}
In conclusion, this thesis demonstrates that the principles of compiler design and functional abstraction can be highly
effective when applied to musical composition.
By bridging the gap between algorithmic logic and the rigid binary specifications of the MIDI protocol, the resulting
domain-specific language empowers composers to explore complex structural arrangements with precision,
confidence, and speed.
